\section*{第 \ref{ch:b2:waves-on-strings} 章——弦与杆上的波：达朗贝尔方程}

\begin{solution}{exo:b2:waves-on-strings:1}
$c = 2Lf = \qty{143}{m/s}$；$T = \mu c^2 = \qty{61}{N}$；弦上 $\lambda = 2L = \qty{1.3}{m}$，
空气中 $340/110 = \qty{3.1}{m}$；第五品：$110 \times 2^{5/12} =
\qty{147}{Hz}$.
\end{solution}

\begin{solution}{exo:b2:waves-on-strings:2}
$\sqrt{E/\rho}$：钢 \qty{5.1}{km/s}，铝 \qty{5.1}{km/s}，铅
\qty{1.2}{km/s}。\qty{1}{km}：沿钢轨需 \qty{0.20}{s}，经空气需 \qty{2.9}{s} ——
先听到的是钢轨传来的声音，两声相隔 \qty{2.7}{s}。
\end{solution}

\begin{solution}{exo:b2:waves-on-strings:3}
$400$、$600$、$800$ Hz 分别对应 $1$、$2$、$3$ 个内部波节；第 3 个模式的波节位于
$L/3$ 和 $2L/3$。$T \times 2$：$200\sqrt2 = \qty{283}{Hz}$；$L/2$：\qty{400}{Hz}；$\mu
\times 3$：$200/\sqrt3 = \qty{115}{Hz}$。
\end{solution}

\begin{solution}{exo:b2:waves-on-strings:4}
$c = \sqrt{45/0.2} = \qty{15}{m/s}$；\qty{24}{m} 所需时间为 \qty{1.6}{s}；脉冲返回时
倒立（$r = -1$）；若为自由端反射，则正立（$r = +1$）。
\end{solution}

\begin{solution}{exo:b2:waves-on-strings:5}
(a) $\mu = \rho\pi d^2/4 = \qty{6.1}{g/m}$；$c = 2Lf = \qty{334}{m/s}$；$T = \mu c^2 =
\qty{690}{N}$。(b) $\sigma = T/A = \rho c^2 = \qty{0.87}{GPa}$，安全裕度 $2.3$。(c) $T/4 =
\qty{170}{N}$，应力相同（仅取决于 $\rho c^2$）。(d) 长弦
不适用此情形；粗的光滑金属丝刚度过大（强烈的
非简谐性），也难以绕过桥面；铜绕组增加了质量，
却不增加刚度。
\end{solution}

\begin{solution}{exo:b2:waves-on-strings:6}
(a) $m_n = 4L^2f^2\mu/n^2g = 6.6/n^2$ kg：$6.6$、$1.65$、$0.73$、\qty{0.41}{kg}。(b)
$n^2 = 13.2$，$n = 3.6$：不共振，位于第 3 与第 4 模式之间（最接近者为 4，
\qty{0.41}{kg}）。(c) 实际上相当于波节：驱动端的振幅
与波腹相比极小。(d) 每个周期内能量同相输入，
并不断累积；阻尼（空气、内摩擦、经滑轮的漏失）
限制了振幅。
\end{solution}

\begin{solution}{exo:b2:waves-on-strings:7}
$c = \qty{100}{m/s}$，$k = \qty{6.3}{rad/m}$，$\lambda = \qty{1.0}{m}$；$v_{\max} = \omega A =
\qty{1.3}{m/s}$，$a_{\max} = \omega^2A = \qty{790}{m/s^2}$；斜率 $kA = 0.013$；$Z =
\sqrt{T\mu} = \qty{1.0}{kg/s}$；$\langle\mathcal P\rangle = \tfrac12Z\omega^2A^2 = \qty{0.79}{W}$；
每个波长内 $\mathcal P/f = \qty{7.9}{mJ}$。
\end{solution}

\begin{solution}{exo:b2:waves-on-strings:8}
(a) $c_1 = \qty{100}{m/s}$，$c_2 = \qty{50}{m/s}$；$Z_1 = \qty{0.20}{kg/s}$，$Z_2 = \qty{0.40}{kg/s}$。
(b) 由轻弦入射：$r = -1/3$，$t = 2/3$；由重弦入射：$r =
+1/3$，$t = 4/3$。(c) 两种情形下均有 $R = 1/9$，$\mathcal T = 4 \times 0.08/0.36 = 8/9$。
(d) 透射波：高 \qty{0.67}{cm}，宽 \qty{5}{cm}（时长相同，
速率减半）；反射波：\qty{-0.33}{cm}，\qty{10}{cm}。
\end{solution}

\begin{solution}{exo:b2:waves-on-strings:9}
(a) $E = F\ell/A\Delta\ell = 80 \times 2/(7.85 \times 10^{-7} \times 10^{-3}) = \qty{2.0e11}{Pa}$。
(b) $c = \qty{5.1}{km/s}$。(c) 纵波：$c/2\ell = \qty{1.3}{kHz}$；横波
$\sqrt{T/\mu} = \qty{114}{m/s}$，$f = \qty{29}{Hz}$ ——张力使金属丝
承受 \qty{100}{MPa} 的应力，相当于 $5 \times 10^{-4}$ 倍的 $E$，因此横波速率为
纵波速率的 $\sqrt{5 \times 10^{-4}}$ 倍。(d) $k = Ea = \qty{50}{N/m}$；$m
= \rho a^3 = \qty{1.2e-25}{kg}$（约 $70$ 个原子质量单位）。
\end{solution}

\begin{solution}{exo:b2:waves-on-strings:10}
(a) $b_n = \frac2L\int_0^Ly(x, 0)\sin(n\pi x/L)\dd x$，其中三角形 $y =
2hx/L$ 位于 $[0, L/2]$ 上（并关于中点对称）：分部积分得 $b_n = 8h\sin(n\pi/2)
/n^2\pi^2$。(b) 偶数 $n$ 对应的系数为零：中点是每个偶数模式的波节，
正是被拨动的点。(c) $E_n = \tfrac14\mu L\omega_n^2b_n^2 \propto n^2/n^4 = 1/n^2$
（奇数 $n$）：基频所占比例为 $1/\sum_{\text{奇数}}n^{-2} = 8/\pi^2 = 81\%$。(d)
$b_n \propto \sin(n\pi/5)/n^2$：$n = 5, 10, \dots$ 时为零。靠近桥的位置，
因子 $\sin(n\pi x_0/L)$ 对低阶模式很小，对高阶模式则持续增大：
高次谐音相对更多，声音更明亮。
\end{solution}

\begin{solution}{exo:b2:waves-on-strings:11}
(a) $-m\omega^2 = k(2\cos ka - 2) = -4k\sin^2(ka/2)$。(b) $v_\varphi = 2\omega_0\sin(ka/2)
/k$，$v_g = \omega_0a\cos(ka/2)$；二者均 $\to \omega_0a = c$（当 $ka \to 0$ 时）；在 $ka = \pi$ 处：
$v_\varphi = 2c/\pi$，$v_g = 0$。(c) 相邻质点反相运动：这是一个驻波
，不输运能量；$k$ 超过 $\pi/a$ 时给出的位移集合与
某个 $k' = k - 2\pi/a$ 的相同 ——并无新的波。(d) $\omega_0 = c/a =
\qty{2.0e13}{rad/s}$，$f_{\max} = \omega_0/\pi = \qty{6.5}{THz}$；对 \qty{1}{MHz}：$\lambda =
\qty{5.1}{mm} = 2 \times 10^7a$，$ka \sim 3 \times 10^{-7}$：完全无色散。
\end{solution}

\begin{solution}{exo:b2:waves-on-strings:12}
(a) $T(x) = \mu gx$，$c = \sqrt{gx}$。(b) $t = \int_0^L\dd x/\sqrt{gx} = 2\sqrt{L/g} =
\qty{2.0}{s}$；脉冲的上部比下部传播得快，因此
脉冲被拉长。(c) $c = \sqrt{T/\mu}$ 向鞭梢方向增大。(d) 能量 $\approx
\mathcal P\tau = ZV^2\tau = \sqrt{T\mu}V^2\tau$ 为常量：$V \propto \mu^{-1/4}$；
$10^4$ 倍的 $\mu$ 差异使速率增至 $\times10$，即 \qty{100}{m/s} ——真实的鞭子还从展开的环中
获得其余能量，这使能量进一步集中，
鞭梢超过了声速。
\end{solution}

\begin{solution}{pb:b2:waves-on-strings:1}
\textbf{1.} $\mu = \rho\pi d^2/4 = \qty{6.1}{g/m}$；$c = 2Lf = \qty{334}{m/s}$；$T = \mu c^2
= \qty{685}{N}$.

\textbf{2.} $\sigma = \rho c^2 = \qty{0.87}{GPa}$：安全裕度 $2.3$。$T \propto f^2$：
升高半音时增加 $12\%$，高出一个八度时增至四倍，
即 \qty{3.5}{GPa} ——金属丝将断裂。

\textbf{3.} $230 \times 685 \approx \qty{160}{kN}$，约十六吨。

\textbf{4.} $c = 2Lf = \qty{105}{m/s}$；在 \qty{685}{N} 下：$\mu = T/c^2 = \qty{63}{g/m}$，
即直径 \qty{3.2}{mm} 的钢杆 ——刚度过大，难以简谐振动，
也难以绕过桥面；以铜绕制的细芯
既提供了质量，又不增加刚度。

\textbf{5.} $440$, $880$, $1320$, $1760$, $2200$, $2640$, $3080$, \qty{3520}{Hz}: $2$,
$4$、$8$ 为八度；$3$、$6$ 为八度之上的纯五度；$5$ 为两个八度之上的
纯大三度；$7$（$3080/1760 = 1.75$）为降
小七度——不协和。

\textbf{6.} $E = \tfrac14\mu L\omega^2A^2 = 0.25 \times 6.1 \times 10^{-3} \times 0.38 \times (2\pi
\times 440)^2 \times (5 \times 10^{-4})^2 = \qty{1.1}{mJ}$.

\textbf{7.} $A\sin kx\cos\omega t = \tfrac12A[\sin(kx - \omega t) + \sin(kx + \omega t)]$：两列
振幅为 \qty{0.25}{mm} 的波；$Z = \mu c = \qty{2.0}{kg/s}$；每列携带的功率为
$\tfrac12Z\omega^2(A/2)^2 = \qty{0.49}{W}$.

\textbf{8.} $\dot y(x, 0) = \sum_nB_n\omega_n\sin(n\pi x/L)$；乘以
$\sin(m\pi x/L)$ 并积分，利用 $\int_0^L\sin\sin = \tfrac L2\delta_{mn}$。

\textbf{9.} 被积函数为 $v_0\sin(n\pi x/L)$，在宽度为 $w$ 的
区间上（围绕 $x_0$），其中正弦近似为常数。

\textbf{10.} $\sin(n\pi/8) = 0$，对 $n = 8, 16, \dots$ 成立；又因 $B_n/B_1 =
\sin(n\pi/8)/(n\sin(\pi/8))$，第七谐音降为基音的 $1/7$，
而协和的 $2$ 至 $6$ 次谐音保持在 $0.9$ 至 $0.3$：
不协和的七度被驯服了。

\textbf{11.} 两个速度脉冲自受击点沿相反
方向传出，在两端倒相反射，穿过后重新合成；
初始状态经 $2L/c = 1/f_1 = \qty{2.3}{ms}$ 后重现。

\textbf{12.} $B = \pi^3 \times 2 \times 10^{11} \times 10^{-12}/(64 \times 685 \times 0.144) =
9.8 \times 10^{-4}$；$\sqrt{1 + 64B} = 1.031$：$+3.1\%$，$1200\log_2(1.031) = 53$ 音分。

\textbf{13.} 低音的谐音比精确的
倍数偏高；为使其与高音的基音重合
（并避免拍音），高音被略微调高。

\textbf{14.} $Z_s = \mu c = \qty{2.0}{kg/s}$; $r = (2 - 2000)/2002 = -0.998$.

\textbf{15.} $1 - r^2 \approx 4Z_s/Z_b = 4 \times 10^{-3}$.

\textbf{16.} 每秒 $c/2L = 440$ 次反射：衰减率 $440 \times 4 \times
10^{-3} = \qty{1.8}{s^{-1}}$，$\tau = \qty{0.56}{s}$；$\ln10^6/1.8 = \qty{7.7}{s}$ 对应
\qty{60}{dB}.

\textbf{17.} $1.1\,\text{mJ} \times 1.8\,\text{s}^{-1} = \qty{2}{mW}$——两百
倍于安静的交谈；它正是你所听到的。

\textbf{18.} 一根 \qty{1}{mm} 粗的弦几乎不推动空气：空气从
它旁边绕过，其辐射功率可忽略不计。音板
通过琴桥接收弦的能量，并带动一平方米的空气
运动：这是一种按面积实现的阻抗匹配。

\textbf{19.} \qty{440}{Hz} 是 A$_3$ 的第二谐音：琴桥以
该弦的某一振动模式驱动它——即感应共振。

\textbf{20.} 提高 $Z_b$：每次反射漏失的能量减少，音
持续更久——但也更弱。

\textbf{21.} $c = \qty{5.1}{km/s}$：\qty{1.0}{s} 在钢轨中，\qty{15}{s} 在空气中。

\textbf{22.} $50/5064 = \qty{10}{ms}$；自由端的应力必须为零：
密部以疏部的形式返回。

\textbf{23.} $\lambda = \qty{2.5}{mm}$；$Z_{\text{钢}} = \rho c = \qty{4e7}{kg.m^{-2}.s^{-1}}$
与 $400$ 相比：$r \approx -1$，发生全反射：任何开口裂纹，无论
多细，都会返回一个回声。

\textbf{24.} $\Delta t = d(1/3.5 - 1/6) = d \times \qty{0.119}{s/km}$: $d = \qty{170}{km}$.

\textbf{25.} 弦 \qty{334}{m/s}（$T$、$\mu$）；钢轨 \qty{5.1}{km/s}（$E$、$\rho$）；
岩石 P 波 \qty{6}{km/s}（压缩模量，$\rho$）；岩石 S 波 \qty{3.5}{km/s}
（剪切模量，$\rho$）；空气 \qty{340}{m/s}（$\gamma P$、$\rho$——下一章）。
\end{solution}
